% Originally: content/week-03.tex, \subsection{Why do we care about compactness?}

\section{Why compactness matters}
\label{sec:why_compactness_matters}

% Sonnet 5 (Medium): added \label on each item so these three facts (heavily cited by name in
% later chapters, e.g. the extreme value theorem) can be \cref-ed instead of just named in prose.
\begin{enumerate}[label=\textbf{(\alph*)}]
  \item \label{item:continuity_preserves_compact}
    \textbf{Preserved by continuous functions:} $X$, $Y$ metric spaces, $K \subseteq X$
    compact, $f : X \to Y$ continuous. Then $f(K) \subseteq Y$ is compact.
  \item \label{item:extreme_value_theorem}
    \textbf{Weierstrass theorem / Extreme value theorem:} $X$ metric space,
    $K \subseteq X$ compact, $f : K \to \mathbb{R}$ continuous. Then $f$ attains its maximum and
    minimum in $K$. In particular, $f$ is bounded. This is proved as
    \cref{lem:sequential_compactness_extreme_values}.
    \begin{ainote}
    The official script's statement of this theorem (Cor.\ 9.79) writes the conclusion as
    \qt{there exist $\bar x \in X$ such that $f(\bar x) = \sup_K f$}, with $\bar x$ ranging over
    the whole ambient space $X$ rather than over $K$; the proof given there repeats the slip,
    writing $f^{-1}(\sup f(X))$ where it must mean $f^{-1}(\sup f(K))$. As stated, the claim is
    false as soon as $f$ is unbounded off $K$: e.g.\ $K:=\{0\}\subset X:=\mathbb{R}$,
    $f(x):=x$. The maximiser $\bar x$ must be sought in $K$, as in the statement above.
    \end{ainote}
  \item \label{item:uniform_continuity_compact}
    \textbf{Uniform continuity:} $X$, $Y$ metric spaces, $f : X \to Y$ continuous. If $X$ is
    compact, then $f$ is uniformly continuous (\cref{def:uniformly_continuous}). This is proved
    as \cref{prop:compact_implies_uniformly_continuous}, via the Lebesgue number of the cover
    $\{f^{-1}(B_{\varepsilon/2}(y))\}_{y}$.
\end{enumerate}

% Sonnet 5 (Medium): illustrates item (b), the extreme value theorem.
\begin{center}
\begin{tikzpicture}[>=Latex, scale=1.2]
  \draw[->] (-0.3,0) -- (3.5,0) node[right] {$x$};
  \draw[->] (0,-0.3) -- (0,2.8) node[above] {$y$};

  \draw[green!60!black, thick] plot[smooth, tension=0.8] coordinates {(0.3,0.8) (0.9,2.3) (2.1,0.4) (3.0,1.8)};

  \draw (0.3,0.05) -- (0.3,-0.05) node[below] {$a$};
  \draw (0.9,0.05) -- (0.9,-0.05) node[below] {$c$};
  \draw (2.1,0.05) -- (2.1,-0.05) node[below] {$d$};
  \draw (3.0,0.05) -- (3.0,-0.05) node[below] {$b$};

  \draw[dashed, red!80!black] (0,2.3) node[left, font=\small] {$f(c)$} -- (0.9,2.3) -- (0.9,0);
  \fill[red!80!black] (0.9,2.3) circle (2pt);

  \draw[dashed, blue!80!black] (0,0.4) node[left, font=\small] {$f(d)$} -- (2.1,0.4) -- (2.1,0);
  \fill[blue!80!black] (2.1,0.4) circle (2pt);

  \node[below, font=\small, align=center] at (1.7,-0.55)
    {\textbf{(b)} on the compact $K=[a,b]$, $f$ attains its maximum at $c$ and its minimum at $d$};
\end{tikzpicture}
\end{center}

Both extrema are attained \emph{inside} $K$. That is the content of \textbf{(b)}, and it is
exactly what fails on a non-compact domain: on the open interval $(a,b)$ the same $f$ may
approach a supremum it never reaches.

% Quelle: Diego Torres Tejeda/Notes - 02.03 - Diego Torres Tejeda.pdf, pp. 8--10
\begin{remark}[What compactness and connectedness are \emph{for}: local to global]
\label{rem:local_to_global}
The three facts above look like a list of unrelated services compactness happens to provide.
Diego Torres Tejeda's notes give the idea that unifies them, and it is worth having early
because it also explains why connectedness sits in the same chapter.

\textbf{Local-to-global arguments} are a recurring theme across mathematics. It is often easy to
check a property \emph{locally}, on some ball $B_\varepsilon(x)$ around each point, while what one
actually wants is the property \emph{globally}, on all of $X$. Compactness and
connectedness are the two hypotheses that license that jump. Schematically, a local-to-global
statement reads:
\begin{quote}
\itshape
If for every $x \in X$ there is $\varepsilon>0$ such that $B_\varepsilon(x)$ has property $P$,
then $X$ has property $P$.
\end{quote}
Each of the two hypotheses powers its own version, and the proofs are three lines each.

\textbf{Connectedness turns \qt{locally constant} into \qt{constant}.} Let $(X,d)$ be connected
and $f : X \to \mathbb{R}$ locally constant. Every level set $f^{-1}(a)$ is then open. Fix
$a \in f(X)$ and split
\[X = \underbrace{f^{-1}(a)}_{=:U} \ \sqcup \underbrace{\bigcup_{b \neq a}f^{-1}(b)}_{=:V},\]
a disjoint union of two open sets. Since $U \neq \emptyset$, connectedness forces $V=\emptyset$,
i.e.\ $f \equiv a$.

\textbf{Compactness turns \qt{locally bounded} into \qt{bounded}.} Let $X$ be compact and
suppose each $x$ has an open $U_x \ni x$ on which $f$ is bounded. The $\{U_x\}_{x\in X}$ cover
$X$, so finitely many suffice, say $U_1,\dots,U_n$, and then
\[\sup_{x \in X}|f(x)| \;=\; \max_{1\leq i\leq n}\ \sup_{x \in U_i}|f(x)| \;<\; \infty,\]
because a \emph{maximum of finitely many} finite numbers is finite. A continuous $f$ is locally
bounded, since $|f|<|f(x)|+1$ near each $x$, so this recovers the boundedness half of
\cref{item:extreme_value_theorem}.

Read this way the two proofs are the same proof twice, with different notions of \qt{finitely
many}: connectedness says a partition into two open pieces must be trivial, compactness says a
cover can be taken finite. Both convert local information into global. Every application in
\Cref{item:continuity_preserves_compact,item:extreme_value_theorem,item:uniform_continuity_compact}
is an instance. Uniform continuity is the most visible of them, since there the whole difficulty
is that the $\delta$ supplied by continuity is local and one needs a single global one
(\cref{rem:continuity_vs_uniform_continuity}).
\end{remark}

% The exercise below overlaps heavily with ex:topological_properties_table above. Both are Sascha
% Brack's, from two different files (his Week 2 Friday and Week 3 Monday notes), and he evidently
% reused the self-test. Both are kept: deciding that one is redundant would be a judgement about
% his notes rather than about this transcription.
% Transition: Claude Opus 5 (High)
The exercise below covers the same ground as \cref{ex:topological_properties_table}, and is worth
doing anyway. It adds the explicit \textbf{?} answers for the cases where the data is
insufficient, a trap in the final item involving two \emph{parallel} spanning vectors, and a
four-column solution table which is where the lesson actually lands. Read
\cref{ex:topological_properties_table} as the warm-up and this one as the real test.

% Supplement: Sascha Brack/Class Notes/Week_02_Notes_Friday_Updated.pdf, p. 8
\begin{exercise}[Classify these sets]
\label{ex:classification_table}
Sascha Brack sets this as a self-test at the end of his compactness lecture, and it is the
fastest way to find out whether the four notions have actually separated in your head. For each
subset of $\mathbb{R}^n$ (standard metric), decide whether it is \textbf{open}, \textbf{closed},
\textbf{compact}, \textbf{complete}. Throughout, $f$ denotes a continuous function and
\textbf{?} means \qt{not decidable from the information given}.
\begin{enumerate}[label=\textbf{\arabic*.}]
  \item $[0,1] \subseteq \mathbb{R}$
  \item $\mathbb{Q} \subseteq \mathbb{R}$
  \item $B_1(0) \subseteq \mathbb{R}^3$
  \item $\left\{\tfrac1n : n \geq 1\right\} \subseteq \mathbb{R}$, and separately
    $\{0\} \cup \left\{\tfrac1n : n \geq 1\right\}$
  \item $\bigcup_{n\geq1}B_{1/n}(n) \subseteq \mathbb{R}$
  \item $(-\infty,1] \subseteq \mathbb{R}$
  \item $f([0,1])$, \quad $f^{-1}([0,1])$, \quad $f^{-1}\big((0,1)\big)$
  \item $S := \operatorname{span}\!\left\{\begin{pmatrix}1\\0\end{pmatrix},
    \begin{pmatrix}3\\0\end{pmatrix}\right\} \subseteq \mathbb{R}^2$
\end{enumerate}
\end{exercise}

% Quelle: Diego Torres Tejeda/Notes - 16.03 - Diego Torres Tejeda.pdf, p. 1 -- Sonnet 5 (Medium)
\begin{remark}[What continuous maps do and do not preserve]
Diego's exercise-class notes collect \cref{item:continuity_preserves_compact} and the analogous
fact for connectedness (\cref{thm:continuity_preserves_connected}) into one table, together with
the dual statement for preimages, and a caution worth stating explicitly:
\begin{center}
\begin{tabular}{lll}
& \textbf{Preserved by $f^{-1}$} & \textbf{Preserved by $f$} \\
\hline
open & yes (\cref{def:continuous}) & not in general \\
closed & yes & not in general \\
compact & not in general & yes (\cref{item:continuity_preserves_compact}) \\
connected & not in general & yes (\cref{thm:continuity_preserves_connected}) \\
\end{tabular}
\end{center}
\textbf{Caution:} completeness and boundedness are \emph{not} topological properties, so in
general they are preserved neither by continuous maps nor by their inverses. For instance,
$\tan : (-\tfrac\pi2,\tfrac\pi2) \to \mathbb{R}$ is a continuous bijection carrying a bounded set
onto an unbounded one.
\end{remark}


% --- exercises relocated here from the dissolved sheet 3 ---

% Originally: content/exercise-sheets/week-03.tex, problem 3.3
% Source: exercises/Ex3_Analysis2_eng.pdf

\begin{exercise}[Open, closed, complete and compact]
\label{ex:3.3}
For each of the following subsets of $\mathbb{R}^N$ say whether they are open / closed / none /
compact (with respect to the standard Euclidean structure). Try to prove your assertions
``efficiently''.
\begin{enumerate}[label=\textbf{\arabic*.}]
  \item $E_1 := \{x \in \mathbb{R}^3 : 0 < x_2 \leq 2\}$
  \item $E_2 := \{x \in \mathbb{R}^n : \sin(|x|) \geq \tfrac{1}{4}\}$
  \item $E_3 := \bigcup_{n \geq 1}\{x \in \mathbb{R}^3 : x_1^4 - \tfrac{1}{n}x_3^2 - x_2^2 > \tfrac{1}{n},\ x_1 + x_2 < 6\}$
  \item $E_4 := \{(x,y) \in (\mathbb{R}^n \times \mathbb{R}^n) : x \cdot y > \tfrac{1}{2}|x||y|\}$
  \item $E_5 := \{X \in \mathbb{R}^{n\times n} : \text{the matrix } X \text{ is invertible}\}$
  \item $E_6 := \{X \in \mathbb{R}^{n\times n} : \text{the matrix } X \text{ is symmetric}\}$
  \item $E_7 := \{X \in \mathbb{R}^{n\times n} : \text{the entries of } X \text{ are either } 0 \text{ or } 1\}$
  \item $E_8 := \{x \in \mathbb{R}^n : |x_i| \leq 6 \text{ for all } 1 \leq i \leq n\}$
  \item $E_9 := E_8 \times E_3 \subset \mathbb{R}^{n+3}$
  \item $E_{10} := E_2 \cap E_8 \subset \mathbb{R}^n$
\end{enumerate}
\exinfo{This exercise is Problem 3.3 of Problem Sheet 3, Analysis II, Spring Semester 2026.}
\end{exercise}

% Originally: content/exercise-sheets/week-03.tex, problem 3.4
% Source: exercises/Ex3_Analysis2_eng.pdf

\begin{exercise}[Multiple choice (open, complete, compact)]
\label{ex:3.4}
Select all the statements below that are true.
\begin{enumerate}[label=\textbf{(\alph*)}]
  \item A nonempty open strict subset of $\mathbb{R}^n$ cannot be compact.
  \item A nonempty open strict subset of $\mathbb{R}^n$ cannot be complete.
  \item A complete subset of $\mathbb{R}^n$ contains all its accumulation points.
  \item Countable union of complete subsets of $\mathbb{R}^n$ is complete.
  \item A subset is closed in $\mathbb{R}^n$ if and only if it is complete as a metric space
    itself (with the distance inherited from $\mathbb{R}^n$).
\end{enumerate}
\exinfo{This exercise is Problem 3.4 of Problem Sheet 3, Analysis II, Spring Semester 2026.}
\end{exercise}

% Source: old_exams/examFS24.pdf (21 August 2024), Wahr/Falsch-Fragen, Aufgabe 1, p. 2
% Extractor: Gemini 3.6 Flash (High)
% Correction: Claude Opus 5 (High) --- \subset -> \subseteq, and part (c) added back: the original
% also asks whether U is *simply* connected, which is the one item on that true/false block with a
% trap in it, and Flash dropped it. Citation corrected: the source is the true/false section, not a
% "Part A".
\begin{exercise}[Continuous image of an annulus]
\label{ex:annulus_continuous_image}
Consider the closed annulus $U := \{(x,y) \in \mathbb{R}^2 \mid 1 \leq x^2+y^2 \leq 4\}$ and the continuous function $f : U \to \mathbb{R}$, $f(x,y) := \sin(xy) - y^4$.
\begin{enumerate}[label=\textbf{(\alph*)}]
  \item Determine whether $U$ is compact and whether $U$ is connected.
  \item Prove that the image $f(U) \subseteq \mathbb{R}$ is a compact, connected interval $[a,b]$ for some real numbers $a < b$.
  \item Is $U$ simply connected?
\end{enumerate}
\exinfo{This exercise is taken from the exam of 21 August 2024, where it appears as the first block
of true/false questions. It is reformulated here as a proof exercise.}
\end{exercise}

% Source: old_exams/fs2023/Prüfung.pdf (9 August 2021), Teil B, Aufgabe 3, p. 15
% Extractor: Claude Opus 5 (High)
\begin{exercise}[Nearest point to the origin]
\label{ex:nearest_point_to_origin}
Let $A \subseteq \mathbb{R}^n$ be non-empty and let $f : A \to \mathbb{R}_{\ge 0}$, $f(x) :=
\lvert x\rvert$, be the restriction of the Euclidean norm to $A$. Prove the following.
\begin{enumerate}[label=\textbf{(\alph*)}]
  \item If $A$ is closed, then $f$ attains a minimum on $A$: there is $P_0 \in A$ with
    $f(P_0) \le f(P)$ for all $P \in A$.
  \item Statement \textbf{(a)} is false in general when $A$ is not assumed closed.
  \item If $A$ is convex, then the minimum is attained at no more than one point of $A$.
  \item If $0 \notin A$, then every minimiser $P_0 \in A$ lies in $\partial A$.
\end{enumerate}
\exinfo{This exercise is taken from the exam of 9 August 2021, Part B, Problem 3.}
\end{exercise}

